\documentclass[10pt,a4paper]{amsart}
\usepackage{amsmath,amssymb,amsthm,graphicx}
\usepackage{verbatim}
\newenvironment{reference}{\comment}{\endcomment}
%\newenvironment{reference}{}{}
\newenvironment{slogan}{\comment}{\endcomment}
\newenvironment{history}{\comment}{\endcomment}

% For commutative diagrams we use Xy-pic
\usepackage[all]{xy}

% We use 2cell for 2-commutative diagrams.
\xyoption{2cell}
% Theorem environments.
%
\theoremstyle{plain}
\newtheorem{theorem}[subsection]{Theorem}
\newtheorem{proposition}[subsection]{Proposition}
\newtheorem{lemma}[subsection]{Lemma}

\theoremstyle{definition}
\newtheorem{definition}[subsection]{Definition}
\newtheorem{example}[subsection]{Example}
\newtheorem{exercise}[subsection]{Exercise}
\newtheorem{situation}[subsection]{Situation}

\theoremstyle{remark}
\newtheorem{remark}[subsection]{Remark}
\newtheorem{remarks}[subsection]{Remarks}

\numberwithin{equation}{subsection}

% Macros
%
\def\lim{\mathop{\mathrm{lim}}\nolimits}
\def\colim{\mathop{\mathrm{colim}}\nolimits}
\def\Spec{\mathop{\mathrm{Spec}}}
\def\Hom{\mathop{\mathrm{Hom}}\nolimits}
\def\Ext{\mathop{\mathrm{Ext}}\nolimits}
\def\SheafHom{\mathop{\mathcal{H}\!\mathit{om}}\nolimits}
\def\SheafExt{\mathop{\mathcal{E}\!\mathit{xt}}\nolimits}
\def\Sch{\mathit{Sch}}
\def\Mor{\mathop{\mathrm{Mor}}\nolimits}
\def\Ob{\mathop{\mathrm{Ob}}\nolimits}
\def\Sh{\mathop{\mathit{Sh}}\nolimits}
\def\NL{\mathop{N\!L}\nolimits}
\def\CH{\mathop{\mathrm{CH}}\nolimits}
\def\proetale{{pro\text{-}\acute{e}tale}}
\def\etale{{\acute{e}tale}}
\def\QCoh{\mathit{QCoh}}
\def\Ker{\mathop{\mathrm{Ker}}}
\def\Im{\mathop{\mathrm{Im}}}
\def\Coker{\mathop{\mathrm{Coker}}}
\def\Coim{\mathop{\mathrm{Coim}}}

% Boxtimes
%
\DeclareMathSymbol{\boxtimes}{\mathbin}{AMSa}{"02}

%
% Macros for moduli stacks/spaces
%
\def\QCohstack{\mathcal{QC}\!\mathit{oh}}
\def\Cohstack{\mathcal{C}\!\mathit{oh}}
\def\Spacesstack{\mathcal{S}\!\mathit{paces}}
\def\Quotfunctor{\mathrm{Quot}}
\def\Hilbfunctor{\mathrm{Hilb}}
\def\Curvesstack{\mathcal{C}\!\mathit{urves}}
\def\Polarizedstack{\mathcal{P}\!\mathit{olarized}}
\def\Complexesstack{\mathcal{C}\!\mathit{omplexes}}
% \Pic is the operator that assigns to X its picard group, usage \Pic(X)
% \Picardstack_{X/B} denotes the Picard stack of X over B
% \Picardfunctor_{X/B} denotes the Picard functor of X over B
\def\Pic{\mathop{\mathrm{Pic}}\nolimits}
\def\Picardstack{\mathcal{P}\!\mathit{ic}}
\def\Picardfunctor{\mathrm{Pic}}
\def\Deformationcategory{\mathcal{D}\!\mathit{ef}}
\usepackage{cleveref}
\crefname{originaldefinition}{original definition}{original definitions}
\Crefname{originaldefinition}{Original definition}{Original definitions}
\makeatletter
\newlabel{groupoids-definition-invariant-open}{{03BC}{source}{}{sourcepointer-1}{}}
\newlabel{groupoids-definition-invariant-open@cref}{{[originaldefinition][0][]03BC}{[1][0][]source}{}{}{}}
\makeatother
\crefname{originallemma}{original lemma}{original lemmas}
\Crefname{originallemma}{Original lemma}{Original lemmas}
\makeatletter
\newlabel{groupoids-lemma-find-invariant-affine}{{03JE}{source}{}{sourcepointer-2}{}}
\newlabel{groupoids-lemma-find-invariant-affine@cref}{{[originallemma][0][]03JE}{[1][0][]source}{}{}{}}
\makeatother
\crefname{originalproposition}{original proposition}{original propositions}
\Crefname{originalproposition}{Original proposition}{Original propositions}
\makeatletter
\newlabel{groupoids-proposition-finite-flat-equivalence}{{03BM}{source}{}{sourcepointer-3}{}}
\newlabel{groupoids-proposition-finite-flat-equivalence@cref}{{[originalproposition][0][]03BM}{[1][0][]source}{}{}{}}
\makeatother
\crefname{originalreference}{original reference}{original references}
\Crefname{originalreference}{Original reference}{Original references}
\makeatletter
\newlabel{lemma-codimension-0-points}{{0BAQ}{source}{}{sourcepointer-4}{}}
\newlabel{lemma-codimension-0-points@cref}{{[originalreference][0][]0BAQ}{[1][0][]source}{}{}{}}
\makeatother
\crefname{originalreference}{original reference}{original references}
\Crefname{originalreference}{Original reference}{Original references}
\makeatletter
\newlabel{lemma-codimension-0-points-dense}{{0BAR}{source}{}{sourcepointer-5}{}}
\newlabel{lemma-codimension-0-points-dense@cref}{{[originalreference][0][]0BAR}{[1][0][]source}{}{}{}}
\makeatother
\crefname{originalreference}{original reference}{original references}
\Crefname{originalreference}{Original reference}{Original references}
\makeatletter
\newlabel{lemma-finite-fibres-presentation}{{03IJ}{source}{}{sourcepointer-6}{}}
\newlabel{lemma-finite-fibres-presentation@cref}{{[originalreference][0][]03IJ}{[1][0][]source}{}{}{}}
\makeatother
\crefname{originalreference}{original reference}{original references}
\Crefname{originalreference}{Original reference}{Original references}
\makeatletter
\newlabel{lemma-subscheme}{{03JH}{source}{}{sourcepointer-7}{}}
\newlabel{lemma-subscheme@cref}{{[originalreference][0][]03JH}{[1][0][]source}{}{}{}}
\makeatother
\crefname{originallemma}{original lemma}{original lemmas}
\Crefname{originallemma}{Original lemma}{Original lemmas}
\makeatletter
\newlabel{properties-lemma-maximal-points-affine}{{01ZV}{source}{}{sourcepointer-8}{}}
\newlabel{properties-lemma-maximal-points-affine@cref}{{[originallemma][0][]01ZV}{[1][0][]source}{}{}{}}
\makeatother
\crefname{originallemma}{original lemma}{original lemmas}
\Crefname{originallemma}{Original lemma}{Original lemmas}
\makeatletter
\newlabel{spaces-lemma-finding-opens}{{02WU}{source}{}{sourcepointer-9}{}}
\newlabel{spaces-lemma-finding-opens@cref}{{[originallemma][0][]02WU}{[1][0][]source}{}{}{}}
\makeatother
\title{An algebraic space with a finite etale cover by a quasi-separated scheme contains a dense open subspace which is a scheme}
\author{The Stacks Project Authors}
\date{stacks-project:spaces-properties.tex:local-snapshot-20261012; GNU FDL 1.2 or later}
\begin{document}
\maketitle
\noindent\textit{Editorial note.} All delivered mathematical text is the author's own, copied from the pinned Stacks Project chapter file spaces-properties.tex; no statement or proof step was rewritten, reordered or removed.\par\medskip
\noindent\textit{Editorial note.} The counted proof is the author's complete argument: the reduction to a point of codimension zero, the construction of an invariant affine open containing the finite fibre of the chart using quasi-separatedness of the cover, and the application of the quotient result for finite flat equivalence relations on affine schemes.\par\medskip
\noindent\textit{Editorial note.} The label lemma-quasi-separated-finite-etale-cover-dense-open-scheme names the counted statement, and the proof environment immediately following it in the source is delivered verbatim as its proof.\par\medskip
\noindent\textit{Editorial note.} The counted statement is carried by the author's own lemma environment in the chapter source.\par\medskip
\noindent\textit{Editorial note.} Pre-open gate: python operations/agent-authoring/tools/precard\_gate.py --source spaces-properties.tex --label lemma-quasi-separated-finite-etale-cover-dense-open-scheme --start 1816 --end 1855 --local-id 28053 returned RESULT=OK (L28053 free, occupancy FREE, public interval clear). Candidate taken from sources/stacks-spaces-properties/candidates-halfA.json.\par\medskip
\noindent\textit{Editorial note.} The statement and proof are copied from the pinned chapter file; the cited inputs are kept as external pointers with their canonical Stacks tags.\par\medskip
\noindent\textit{Allowed standard prerequisites (editorial).}\begin{itemize}
\item Existence of a maximal open subspace which is a scheme. Same chapter, the lemma producing the maximal open subspace of an algebraic space which is a scheme, together with the lemma that points of codimension zero are dense. Supplies the maximal open subspace X' and the reduction of the claim to showing that a point of codimension zero lies in it.
\item Finiteness of the fibres of the chart and codimension of their points. Same chapter, the lemma that the fibres of the map on points from a finite chart are finite, and the lemma identifying the points of codimension zero in a finite fibre as generic points. Supplies the finitely many points eta\_i over x, each a generic point of an irreducible component of U.
\item Finding an invariant affine open containing finitely many points. Properties, the lemma that finitely many points of a quasi-separated scheme lie in an affine open, and Groupoids, the lemma producing an invariant affine open neighbourhood. Supplies the invariant affine open W containing the eta\_i.
\item The quotient of a finite flat equivalence relation on an affine scheme is a scheme. Groupoids, the proposition computing the quotient of a finite flat equivalence relation as a scheme, and Spaces, the lemma identifying the quotient of an invariant open as an open subspace. Supply the conclusion that W/R\_W is a scheme and an open subspace of X containing x.
\end{itemize}\medskip
\ifdefined\hypertarget\hypertarget{sourcepointer-1}{}\fi
\noindent\textit{Source reference.} Original definition 03BC: sources/stacks-chow/tags.txt, source lines 4343--4343. Prerequisite pointer invoked at proof line(s) 1847; it names a standard fact of a sibling chapter that the counted argument uses, and every step of the counted proof built on it is retained verbatim in the delivered body.\par\medskip
\ifdefined\hypertarget\hypertarget{sourcepointer-2}{}\fi
\noindent\textit{Source reference.} Original lemma 03JE: sources/stacks-chow/tags.txt, source lines 4636--4636. Prerequisite pointer invoked at proof line(s) 1843; it names a standard fact of a sibling chapter that the counted argument uses, and every step of the counted proof built on it is retained verbatim in the delivered body.\par\medskip
\ifdefined\hypertarget\hypertarget{sourcepointer-3}{}\fi
\noindent\textit{Source reference.} Original proposition 03BM: sources/stacks-chow/tags.txt, source lines 4353--4353. Prerequisite pointer invoked at proof line(s) 1851; it names a standard fact of a sibling chapter that the counted argument uses, and every step of the counted proof built on it is retained verbatim in the delivered body.\par\medskip
\ifdefined\hypertarget\hypertarget{sourcepointer-4}{}\fi
\noindent\textit{Source reference.} Original reference 0BAQ: sources/stacks-spaces-properties/pinned-original-source.tex, source lines 1554--1554. Prerequisite pointer invoked at proof line(s) 1837; it names a standard fact of the same chapter that the counted argument uses, and every step of the counted proof built on it is retained verbatim in the delivered body.\par\medskip
\ifdefined\hypertarget\hypertarget{sourcepointer-5}{}\fi
\noindent\textit{Source reference.} Original reference 0BAR: sources/stacks-spaces-properties/pinned-original-source.tex, source lines 1578--1578. Prerequisite pointer invoked at proof line(s) 1830; it names a standard fact of the same chapter that the counted argument uses, and every step of the counted proof built on it is retained verbatim in the delivered body.\par\medskip
\ifdefined\hypertarget\hypertarget{sourcepointer-6}{}\fi
\noindent\textit{Source reference.} Original reference 03IJ: sources/stacks-spaces-properties/pinned-original-source.tex, source lines 1155--1155. Prerequisite pointer invoked at proof line(s) 1835; it names a standard fact of the same chapter that the counted argument uses, and every step of the counted proof built on it is retained verbatim in the delivered body.\par\medskip
\ifdefined\hypertarget\hypertarget{sourcepointer-7}{}\fi
\noindent\textit{Source reference.} Original reference 03JH: sources/stacks-spaces-properties/pinned-original-source.tex, source lines 1782--1782. Prerequisite pointer invoked at proof line(s) 1828; it names a standard fact of the same chapter that the counted argument uses, and every step of the counted proof built on it is retained verbatim in the delivered body.\par\medskip
\ifdefined\hypertarget\hypertarget{sourcepointer-8}{}\fi
\noindent\textit{Source reference.} Original lemma 01ZV: sources/stacks-chow/tags.txt, source lines 2610--2610. Prerequisite pointer invoked at proof line(s) 1839; it names a standard fact of a sibling chapter that the counted argument uses, and every step of the counted proof built on it is retained verbatim in the delivered body.\par\medskip
\ifdefined\hypertarget\hypertarget{sourcepointer-9}{}\fi
\noindent\textit{Source reference.} Original lemma 02WU: sources/stacks-chow/tags.txt, source lines 3817--3817. Prerequisite pointer invoked at proof line(s) 1853; it names a standard fact of a sibling chapter that the counted argument uses, and every step of the counted proof built on it is retained verbatim in the delivered body.\par\medskip
\begin{lemma}
\label{lemma-quasi-separated-finite-etale-cover-dense-open-scheme}
Let $S$ be a scheme. Let $X$ be an algebraic space over $S$.
If there exists a finite, \'etale, surjective morphism
$U \to X$ where $U$ is a quasi-separated scheme, then
there exists a dense open subspace $X'$ of $X$ which is a scheme.
More precisely, every point $x \in |X|$ of codimension $0$ in $X$
is contained in $X'$.
\end{lemma}

\begin{proof}
Let $X' \subset X$ be the maximal open subspace which is a scheme
(Lemma \ref{lemma-subscheme}).
Let $x \in |X|$ be a point of codimension $0$ on $X$.
By Lemma \ref{lemma-codimension-0-points-dense}
it suffices to show $x \in X'$.
Let $U \to X$ be as in the statement of the lemma.
Write $R = U \times_X U$ and denote $s, t : R \to U$ the projections as usual.
Note that $s, t$ are surjective, finite and \'etale.
By Lemma \ref{lemma-finite-fibres-presentation}
the fibre of $|U| \to |X|$ over $x$ is finite, say
$\{\eta_1, \ldots, \eta_n\}$. By Lemma \ref{lemma-codimension-0-points}
each $\eta_i$ is the generic point of an irreducible component of $U$.
By Properties, Lemma \ref{properties-lemma-maximal-points-affine}
we can find an affine open $W \subset U$ containing
$\{\eta_1, \ldots, \eta_n\}$
(this is where we use that $U$ is quasi-separated). By
Groupoids, Lemma \ref{groupoids-lemma-find-invariant-affine}
we may assume that $W$ is $R$-invariant.
Since $W \subset U$ is an $R$-invariant affine open, the restriction
$R_W$ of $R$ to $W$ equals $R_W = s^{-1}(W) = t^{-1}(W)$ (see
Groupoids, Definition \ref{groupoids-definition-invariant-open}
and discussion following it). In particular the maps $R_W \to W$ are
finite \'etale also. It follows that $R_W$ is affine.
Thus we see that $W/R_W$ is a scheme, by
Groupoids, Proposition \ref{groupoids-proposition-finite-flat-equivalence}.
On the other hand, $W/R_W$ is an open subspace of $X$ by
Spaces, Lemma \ref{spaces-lemma-finding-opens} and it contains
$x$ by construction.
\end{proof}

\end{document}
